% Shared hyperlink and cross-reference discipline for the Book and its chapters.
% Load this LAST in every preamble: cleveref must follow hyperref and amsthm.
% hyperref itself is loaded earlier with [backref=page] (a load-time option).
% Twin copy under whitepaper/figures/ (byte-identical; checked by
% scripts/generate-mega-whitepaper.test.mjs).
%
% What a reader gets: live links on chapter titles, section numbers, theorems,
% figures, tables, equations, page numbers, and citations; PDF bookmarks three
% levels deep; and, at the end of every bibliography entry, the pages that cite it.
% Links are quiet: a warm dark ink one step lighter than body text, so the
% four part hues keep their meaning and an inline chapter title never competes
% with the part it sits in. Chrome pages (the front-matter map, part openers)
% recolor their own links locally.
\hypersetup{
  colorlinks=true,
  linkcolor=pdinkmuted,
  citecolor=pdinkmuted,
  urlcolor=pdinkmuted,
  filecolor=pdinkmuted,
  anchorcolor=pdink,
  bookmarksnumbered=true,
  bookmarksdepth=3,
  bookmarksopen=true,
  bookmarksopenlevel=1,
  breaklinks=true,
  linktoc=all,
  pdfauthor={Erich Owens},
  pdftitle={\pdchaptertitle},
  pdfsubject={\pdeditiontitle: \pdeditionsubtitle},
}

% Title-based chapter references. \pdchaplink{prefix}{text} is a live link to
% the chapter inside the Book (the Book's preamble defines it before loading this
% file); in a standalone chapter it is plain text. Prose uses \pdchapref, which
% sets the title in italics; the locator map uses \pdchaplink directly.
\providecommand{\pdchaplink}[2]{#2}
\DeclareRobustCommand{\pdchapref}[2]{\pdchaplink{#1}{\emph{#2}}}

% \ifpdbook -- true in the assembled Book, false in a standalone paper.
%
% WHY THIS EXISTS. A drawing that two chapters both \input prints twice in the
% Book, under two figure numbers. The fix is for the chapter that develops the
% idea to keep the figure and the other to point at it -- but a standalone
% paper has nowhere to point, and rewriting the prose outright removed five
% figures from harbor-economy-whitepaper.pdf and left it telling a conference
% reader that a ceremony is "drawn in The Federated Harbor", a document that
% reader does not have.
%
% So the divergence is per-build. A CONDITIONAL and not a two-argument macro:
% the generator inlines each \input, so a figure branch is a whole table or
% tikzpicture, and as a macro argument that is scanned as one token list --
% which is how \pdinbook produced "runaway argument" on the first attempt.
% \ifpdbook skips the untaken branch at the TeX level instead.
%
% Declared unconditionally and defaulting FALSE, which is the standalone
% answer. The Book's preamble sets \pdbooktrue immediately after loading
% this file. A guard here would have to mention \ifpdbook or \newif inside a
% branch TeX may skip, and a conditional token in skipped text is counted as
% a nested \if -- which is how two attempts at a guard both produced
% "Incomplete \ifdefined". Ordering is cheaper than cleverness.
\newif\ifpdbook

% cleveref: "Theorem 3.1", "Figure 2.4", "Section 1.3", capitalised, never
% abbreviated, with the whole phrase clickable. Names are declared explicitly
% because the theorem environments are created before this file is loaded.
\usepackage[nameinlink,capitalise,noabbrev]{cleveref}
\crefname{pdchapter}{Chapter}{Chapters}
\crefname{part}{Part}{Parts}
\crefname{appendix}{Appendix}{Appendices}
\crefname{definition}{Definition}{Definitions}
\crefname{theorem}{Theorem}{Theorems}
\crefname{lemma}{Lemma}{Lemmas}
\crefname{proposition}{Proposition}{Propositions}
\crefname{corollary}{Corollary}{Corollaries}
\crefname{conjecture}{Conjecture}{Conjectures}
\crefname{property}{Property}{Properties}
\crefname{protocol}{Protocol}{Protocols}
\crefname{lsprotocol}{Protocol}{Protocols}
\crefname{stpprotocol}{Protocol}{Protocols}
\crefname{heprotocol}{Protocol}{Protocols}
\crefname{remark}{Remark}{Remarks}
\crefname{lstlisting}{Listing}{Listings}
\crefname{exercise}{Exercise}{Exercises}
\crefname{solution}{Solution}{Solutions}
\Crefname{pdchapter}{Chapter}{Chapters}
\Crefname{lstlisting}{Listing}{Listings}
\Crefname{exercise}{Exercise}{Exercises}
\Crefname{solution}{Solution}{Solutions}

% Back-references: "^ Cited on pages 12, 40, and 87." after every bibliography
% entry, each page number a live link back to the page that cited the work.
% Distinct pages only (#3/#4).
%
% The page numbers used to be wrapped in \color{hhgray} along with the words,
% which painted over the link colour hyperref had just given them: the links
% were there -- 29 annotations on a typical bibliography page -- and looked
% exactly like the grey prose around them, so nobody knew they could click.
% Restoring the quiet link colour would not have fixed it either; pdinkmuted
% against hhgray is 1.53:1, which is no signal at footnote size.
%
% So the numbers take pdcobalt (7.55:1 on the cream, and a hue a reader cannot
% mistake for grey), and the line opens with the return arrow this book already
% uses to send a solution back to its exercise. This is a deliberate exception
% to the quiet-link discipline above, and the reason is that back-matter
% navigation is the one place where the link IS the content: in prose the
% surrounding sentence tells you a phrase is a destination, and in a
% bibliography nothing does.
\renewcommand*{\backref}[1]{}
\renewcommand*{\backrefalt}[4]{%
  \ifcase #3\relax
  \or {\footnotesize{\color{hhgray}$\uparrow$~Cited on page}~{\hypersetup{linkcolor=pdcobalt}#4}{\color{hhgray}.}}%
  \else {\footnotesize{\color{hhgray}$\uparrow$~Cited on pages}~{\hypersetup{linkcolor=pdcobalt}#4}{\color{hhgray}.}}%
  \fi}
\renewcommand*{\backrefsep}{, }
\renewcommand*{\backreftwosep}{ and }
\renewcommand*{\backreflastsep}{, and }
